\hypertarget{the-heart-of-the-machine-ceval.c-and-the-interpreter-loop}{%
\chapter{\texorpdfstring{The Heart of the Machine: \texttt{ceval.c} and
the Interpreter
Loop}{The Heart of the Machine: ceval.c and the Interpreter Loop}}\label{the-heart-of-the-machine-ceval.c-and-the-interpreter-loop}}

To understand Python execution is to understand the main evaluation
loop. In CPython, this resides in
\passthrough{\lstinline!Python/ceval.c!}, specifically in the function
\passthrough{\lstinline!\_PyEval\_EvalFrameDefault!}.

\hypertarget{the-mega-switch-statement}{%
\subsection{69.1 The Mega-Switch
Statement}\label{the-mega-switch-statement}}

Historically, the Python interpreter loop was a giant C
\passthrough{\lstinline!switch!} statement inside a
\passthrough{\lstinline!while!} loop.

\begin{lstlisting}[language=C]
for (;;) {
    opcode = NEXTOPARG();
    switch (opcode) {
        case TARGET(LOAD_FAST):
            // ... load local variable ...
            FAST_DISPATCH();
        case TARGET(BINARY_ADD):
            // ... add two objects ...
            FAST_DISPATCH();
    }
}
\end{lstlisting}

\hypertarget{computed-gotos}{%
\subsubsection{1. Computed Gotos}\label{computed-gotos}}

On compilers that support it (like GCC and Clang), CPython uses
``Computed Gotos.'' Instead of a switch statement (which requires a jump
table lookup and a bounds check for every instruction), it uses a table
of memory addresses. At the end of each opcode's C code, it jumps
directly to the address of the next opcode. This reduces CPU branch
mispredictions and significantly improves performance.

\hypertarget{the-evaluation-stack}{%
\subsection{69.2 The Evaluation Stack}\label{the-evaluation-stack}}

Python is a stack-based VM. * \textbf{The Stack Pointer}:
\passthrough{\lstinline!stack\_pointer!} in C. * \textbf{PUSH/POP}:
These are simple pointer increments/decrements in C. * \textbf{Value
Stack}: An array of \passthrough{\lstinline!PyObject *!}. When you add
two numbers, the pointers to the numbers are popped, the addition is
performed, and the pointer to the new result object is pushed.

\hypertarget{handling-interrupts-and-the-gil}{%
\subsection{69.3 Handling Interrupts and the
GIL}\label{handling-interrupts-and-the-gil}}

The interpreter loop is not just for math; it is the system's heartbeat.
* \textbf{Signal Checking}: Every \(N\) instructions (the ``check
interval''), the loop checks if a signal has arrived from the OS. *
\textbf{Thread Switching}: This is also where the GIL is released and
re-acquired, allowing other threads to take their turn in the VM.

\begin{center}\rule{0.5\linewidth}{0.5pt}\end{center}
